\chapter{Extension: a known common trend}\label{ch:trend}
% Not in the thesis. Lean: Credibility/Trend.lean (notes/time-trends.md §1, steps 1–2)

The world has a common claim trend $g > -1$ and noise that grows with the
level (constant coefficient of variation):
$X_{jt} = (1+g)^t(\mu_j + \sigma_{jt}\varepsilon_{jt})$. Restated to the
target year $n$ (usually $T+1$),
$\tilde X_{jt} = (1+g)^{n-t}X_{jt} = (1+g)^n(\mu_j + \sigma_{jt}\varepsilon_{jt})$:
the model of Chapter~\ref{ch:unequal} in a unit $(1+g)^n$ times smaller.
So nothing new is needed; this chapter records that the fit does not move.

\begin{lemma}[Scaling $K$]\label{lem:glsWeight_smul}
  \lean{Credibility.inv_smul_of_ne_zero, Credibility.invTotal_smul,
        Credibility.glsWeight_smul}
  \leanok
  \uses{def:glsWeight}
  For $c \ne 0$: $(cK)^{-1} = c^{-1}K^{-1}$ (also for singular $K$),
  $(cK)^{-1}_{**} = c^{-1}K^{-1}_{**}$, and the GLS weights of $cK$ are those of $K$.
\end{lemma}
\begin{proof}
  \leanok
  The factor $c^{-1}$ cancels in $b = K^{-1}\bu/K^{-1}_{**}$.
\end{proof}

\begin{definition}[Scaling the blend moments]\label{def:BM_smul}
  \lean{Credibility.BlendMoments.smul}
  \leanok
  \uses{def:errVar}
  Multiply all six second moments of the blend by $c$.
\end{definition}

\begin{lemma}[$z^\ast$ is unit-free]\label{lem:BM_smul}
  \lean{Credibility.BlendMoments.errVar_smul, Credibility.BlendMoments.gap_smul,
        Credibility.BlendMoments.bestZ_smul}
  \leanok
  \uses{def:BM_smul, def:bestZ}
  The error variance and the curvature scale by $c$; for $c \ne 0$, $z^\ast$ is unchanged.
\end{lemma}
\begin{proof}
  \leanok
  $z^\ast$ is a ratio of two expressions linear in the moments.
\end{proof}

\begin{lemma}[Scaling $K$, $\tau^2$ and $v$]\label{lem:Core_moments_smul}
  \lean{Credibility.Core.moments_smul, Credibility.Core.bestZ_smul}
  \leanok
  \uses{def:Core.moments, def:BM_smul}
  For $c \ne 0$ the blend moments of $(cK, c\tau^2, cv)$ are those of
  $(K, \tau^2, v)$ times $c$, so $z_j$ is unchanged.
\end{lemma}
\begin{proof}
  \leanok
  \uses{lem:glsWeight_smul, lem:BM_smul}
  $b$ is unchanged and $s^2 = 1/K^{-1}_{**}$ scales by $c$
  (Lemma~\ref{lem:glsWeight_smul}); then Lemma~\ref{lem:BM_smul}.
\end{proof}

\begin{definition}[Scaled and restated model]\label{def:U_scale}
  \lean{Credibility.Unequal.scale, Credibility.Unequal.detrend}
  \leanok
  \uses{def:Unequal}
  For $c > 0$, $\mathrm{scale}\,c$ replaces $\tau^2$ by $c^2\tau^2$ and every
  $\sigma_{jt}$ by $c\sigma_{jt}$, keeping $\rho$ and the periods $T_j$. The data
  restated to year $n$ with a known trend $g > -1$ have the model
  $\mathrm{detrend}\,g\,n = \mathrm{scale}\,(1+g)^n$.
\end{definition}

\begin{theorem}[Scaling moves no weight]\label{thm:U_scale}
  \lean{Credibility.Unequal.a_scale, Credibility.Unequal.σ2_scale,
        Credibility.Unequal.K_scale, Credibility.Unequal.b_scale,
        Credibility.Unequal.s2_scale, Credibility.Unequal.moments_scale,
        Credibility.Unequal.bestZ_scale, Credibility.Unequal.K_posDef_scale}
  \leanok
  \uses{def:U_scale, thm:U_bestZ_eq}
  In the scaled model the cell weights $a_{jt}$, the GLS weights $b_j$ and
  $z_j$ are unchanged; $\sigma^2_{ij}$, $K$, $s^2$ and all blend moments are
  multiplied by $c^2$, and $K$ stays positive definite.
\end{theorem}
\begin{proof}
  \leanok
  \uses{lem:U_K_apply, lem:glsWeight_smul, lem:Core_moments_smul}
  $c^2$ cancels in \eqth{23}, so $\sigma^2_{ij}$ scales by $c^2$, and so does
  $K_{ij} = \sigma^2_{ij}\rho_{ij} + \tau^2\delta_{ij}$
  (Lemma~\ref{lem:U_K_apply}). Then Lemmas~\ref{lem:glsWeight_smul}
  and~\ref{lem:Core_moments_smul}.
\end{proof}

\begin{theorem}[Trend first, then blend]\label{thm:U_detrend}
  \lean{Credibility.Unequal.a_detrend, Credibility.Unequal.K_detrend,
        Credibility.Unequal.b_detrend, Credibility.Unequal.s2_detrend,
        Credibility.Unequal.bestZ_detrend, Credibility.Unequal.K_posDef_detrend}
  \leanok
  \uses{def:U_scale}
  The restated data have $\tilde K = (1+g)^{2n}K$ and
  $\tilde s^2 = (1+g)^{2n}s^2$; $a_{jt}$, $b_j$ and $z_j$ do not depend on $g$ or $n$.
\end{theorem}
\begin{proof}
  \leanok
  \uses{thm:U_scale}
  Theorem~\ref{thm:U_scale} with $c = (1+g)^n$.
\end{proof}

\begin{definition}[The same noise in the scaled model]\label{def:U_Noise_scale}
  \lean{Credibility.Unequal.Noise.scale, Credibility.Unequal.Noise.detrend}
  \leanok
  \uses{def:U_Noise, def:U_scale}
  Moment assumptions for $M$ (Definition~\ref{def:U_Noise}) are moment
  assumptions for the scaled model, with the same $\eta$ and $\varepsilon$.
\end{definition}

\begin{theorem}[Restated data are the thesis model]\label{thm:U_detrend_obs}
  \lean{Credibility.Unequal.Noise.obsVal_scale, Credibility.Unequal.Noise.cohortEst_scale,
        Credibility.Unequal.Noise.portEst_scale, Credibility.Unequal.Noise.blend_scale,
        Credibility.Unequal.Noise.detrend_obs}
  \leanok
  \uses{def:U_Noise_scale}
  With mean $c\mu$, every observation, cohort estimate, portfolio estimate
  and blend of the scaled model is $c$ times that of $M$, outcome by outcome.
  In particular, for any numbering of the years, the trended observation
  restated to year $n$, $(1+g)^{n-t}\,(1+g)^t\mu_{jt}$, is the observation of
  $\mathrm{detrend}\,g\,n$ with mean $(1+g)^n\mu$.
\end{theorem}
\begin{proof}
  \leanok
  \uses{thm:U_scale}
  $\sqrt{c^2\tau^2} = c\tau$; the weights $a$, $b$ are unchanged
  (Theorem~\ref{thm:U_scale}) and the estimates are linear in the data.
\end{proof}

\begin{theorem}[Relativities are trend-free]\label{thm:U_relativity}
  \lean{Credibility.Unequal.Noise.relativity_detrend}
  \leanok
  \uses{thm:U_detrend_obs, thm:U_detrend}
  In the world with trend $g$, the fitted relativity
  $\bigl(z_j\hmu_j + (1-z_j)\hmu\bigr)/\hmu$ of the data restated with $g$, with
  the best $z_j$, equals that of the untrended world for the same outcome: it
  does not depend on $g$ or on the target year. The trend moves only the level.
\end{theorem}
\begin{proof}
  \leanok
  \uses{thm:U_detrend_obs, thm:U_detrend}
  $z_j$ is unchanged (Theorem~\ref{thm:U_detrend}), and numerator and
  denominator are both multiplied by $(1+g)^n$ (Theorem~\ref{thm:U_detrend_obs}).
\end{proof}

\begin{remark}[Finding 21: an estimated trend does move relativities]\label{rem:trend_estimated}
  \uses{thm:U_relativity, thm:U_detrend}
  Theorem~\ref{thm:U_relativity} restates with the \emph{true} $g$. Restating
  with an estimate $\hat g \ne g$ leaves a residual trend
  $((1+g)/(1+\hat g))^t$ that varies across years, so it reweights old against
  recent data within each cohort, and relativities change.
  Counterexample: two cohorts, two years, equal $\sigma$, $\rho = 0$, data
  $(2, 1)$ and $(1, 2)$. Restated with $r = 1+\hat g$, $\hmu_1/\hmu_2 = (2r+1)/(r+2)$:
  $1$ at $\hat g = 0$, $1.032$ at $\hat g = 10\%$. What does hold for any
  $\hat g$: $a$, $b$ and $z$ do not depend on it (Theorem~\ref{thm:U_detrend}),
  and the target year never matters. The error of $\hat g$ is a level shift
  only when every cohort's data have the same time profile.
\end{remark}

% ---------------------------------------------------------------------------
